\documentclass[11pt]{article}
\usepackage[margin=1.15in]{geometry}
\usepackage{amsmath,amssymb,amsthm}
\usepackage{booktabs}
\usepackage[round]{natbib}
\usepackage[hidelinks]{hyperref}
% run-in headings italic, not bold (house style)
\makeatletter
\renewcommand\paragraph{\@startsection{paragraph}{4}{\z@}%
  {3.25ex \@plus1ex \@minus.2ex}{-1em}%
  {\normalfont\normalsize\itshape}}
\makeatother
\newtheorem{theorem}{Theorem}
\newtheorem{proposition}{Proposition}
\newtheorem{remark}{Remark}
\newcommand{\R}{\mathbb{R}}
\newcommand{\E}{\mathbb{E}}

\title{Scalable Inversion of Contests with Correlated Performances,\\
Including Softmax and Multinomial Probit}
\author{Peter Cotton\thanks{\texttt{peter.cotton@microprediction.com}. The reference
implementation (Python), its parity-locked ports (R, Rust, JavaScript) and the seeded
scripts behind every number in this paper live at
\texttt{github.com/microprediction/winning}; this revision's tables were
produced at repository tag \texttt{paper-r3}, package version
\texttt{1.3.0}.}}
\date{August 2026}

\begin{document}
\maketitle

\begin{abstract}
Multinomial probit choice probabilities over $n$ alternatives are
Gaussian orthant integrals, computed by simulation for thirty years,
one expensive integral per alternative. We observe that the choice
event is one-dimensional: whoever wins, wins at some value, so a
single shared survival field prices all $n$ alternatives at once. For
the covariances applied work actually uses (factor, block, nested,
hierarchical: conditional independence given a few shared variables),
the field yields all $n$ probabilities in $O(nLQ)$ operations over a
lattice of $L$ winning values and $Q$ factor quadrature nodes, and
inverts them back to utilities under a fixed, known covariance. For the
factor grammar the same pass gives matrix-free Jacobian--vector
products and any specified single-removal counterfactual at the same
order; the complete removal table has quadratic output size. The Jacobian is a weighted graph Laplacian whose edge
weights are photo-finish tie densities: the geometry of substitution
between alternatives, from the same pass that prices them. The cleanest
evidence is internal: the identical factor conditioning computed one
alternative at a time, on the same nodes and lattice, agrees to
$3\times10^{-15}$ and costs six hundred times more. Against the field
standard, the Geweke--Hajivassiliou--Keane (GHK) simulator at $n=200$
is two hundred times slower at $10^{4}$ draws, with the lattice's
error below the comparison's resolution. Forward
pricing is benchmarked to ten million alternatives, inversion to one
million. Softmax, Thurstone--Mosteller, factor probit, and team and
sector hierarchies are the same two calls, and the same objects serve
selection and routing among correlated populations of models or agents.
\end{abstract}

\section{Introduction}\label{sec:intro}

A decision maker faces $n$ alternatives and chooses the one with the highest
utility. Write $U_i = \mu_i + \varepsilon_i$ with
$\varepsilon\sim N(0,\Sigma)$; the probability that alternative $i$ is chosen is
\begin{equation}\label{eq:orthant}
p_i \;=\; \Pr\{U_i \ge U_j \ \text{for all } j\}
      \;=\; \Pr\{\tilde U \ge 0\},
\qquad \tilde U_j = U_i - U_j,
\end{equation}
an $(n-1)$-dimensional Gaussian orthant integral. This is the multinomial probit
model. Its appeal has been understood since \citet{thurstone1927}: unlike the
logit family it imposes no independence of irrelevant alternatives, because
$\Sigma$ lets alternatives be close substitutes. Its curse has been understood
just as long. The integral \eqref{eq:orthant} has no closed form, and estimation
needs it, and its derivatives, at every trial value of the parameters.

The workhorse answer is the GHK simulator of \citet{geweke1989},
\citet{borschsupan1993}, \citet{keane1994} and \citet{hajivassiliou1996}. GHK
factors $\Sigma$ by Cholesky and writes the orthant probability as a nested
sequence of one-dimensional truncated normal probabilities, each conditioning on
draws for the previous coordinates. Averaging $R$ such importance weights gives an
unbiased, smooth estimate of one $p_i$ at cost $O(n^2)$ per draw. The estimator
made maximum simulated likelihood practical and it remains the default in
textbooks \citep{train2009}. Its costs are also standard: error shrinks as
$O(R^{-1/2})$; the log of a simulated probability is biased at $O(1/R)$;
gradients inherit the simulation noise; and a full probability vector costs
$n$ separate runs. A single likelihood contribution needs only the chosen
alternative's probability, so the $n$-fold cost falls not on the individual
log-likelihood but on everything else estimation touches: share inversion,
aggregate-share and market models, prediction and welfare across the menu,
and the derivative of any of these with respect to all $n$ utilities.

This paper takes a different route to the same object, in a min-wins
convention fixed here once: performances are times, $X_i=-U_i$, and the
lowest time wins. The choice event is one-dimensional. Contestant $i$ wins
when its time beats the field, so
\begin{equation}\label{eq:field-intro}
p_i \;=\; \int f_i(x)\prod_{j\ne i}S_j(x)\,dx,
\end{equation}
where $f_i$ is the density of $X_i$ and $S_j$ the survival of $X_j$:
integrate over the winning time $x$, requiring every rival to come in
slower. For
independent performances \citet{cotton2021} computed \eqref{eq:field-intro} for
all $n$ on a lattice by building the field $\sum_j \log S_j$ once and dividing
each alternative out, and inverted the map from probabilities back to $\mu$. The
present paper extends the field construction to correlated performances by
conditioning: whenever $\Sigma$ makes performances independent given a few
shared variables, the field factorizes and all $n$ probabilities cost
$O(nLQ)$ for a lattice of $L$ points and $Q$ quadrature nodes. For the
factor grammar the derivatives ride the same pass at the same cost:
Jacobian--vector products and each inversion iteration are $O(nLQ)$, with
the Jacobian a weighted graph Laplacian applied as an operator and never
materialized. The block and nested kernels materialize their Jacobians,
and the tree's cross-cluster block is approximate, so the matrix-free
claim is the factor grammar's, not every grammar's. No per-alternative
repetition; error controlled by quadrature,
deterministic for the Gauss--Hermite rules and reproducibly randomized for
scrambled-Sobol nodes.

Only utility contrasts and a normalized scale are identified in multinomial
probit, so the raw covariance is overparameterized, and the contrast
covariance may in principle carry unrestricted parameters. Applied work
nevertheless imposes structured covariance as a matter of course, for
parsimony and for computation \citep{train2009}; the grammars here cover the
structured families in common use. The table lists them: the first six rows
are nested exactly and each is priced and inverted by the same two calls; the
last row is a fitted approximation developed in \S\ref{sec:general}.

\begin{center}\small
\begin{tabular}{lll}
\toprule
model & covariance grammar & source \\
\midrule
Luce / softmax & independent, minimum-Gumbel base, $D=\pi^2/6$ & exact identity \\
Thurstone--Mosteller & independent, normal base & \citet{cotton2021} \\
factor multinomial probit & $VV'+\mathrm{diag}(D)$ & \S\ref{sec:factor} \\
teams, conferences & blocks: rank-$r$ effect per cluster & \S\ref{sec:blocks} \\
market plus sectors & nested: blocks $\times$ global factor & \S\ref{sec:blocks} \\
hierarchies & tree: uniform shared effects & \S\ref{sec:tree} \\
dense $\Sigma$ & fitted grammar $+$ residual factors (approx.) & \S\ref{sec:general} \\
\bottomrule
\end{tabular}
\end{center}

The second motivation is scale. Measured, on one laptop: all-$n$
forward pricing at
$n=10^{4},10^{5},10^{6}$ in $0.18$, $2.7$ and $29$ seconds (rank-one
factor, Rust kernels), a ten-million-contestant block field
in $245$ seconds, and inversion at $n=10^{6}$ in $80$ seconds
(independent) and $22$ minutes (rank-one factor). GHK's measured
complete-vector cost is superquadratic and steepening (log-log slope
$2.15$ from $n=10$ to $200$, $2.8$ from $200$ to $1000$; timings at
$R=10^{3}$ of $0.57$, $5.66$ and $51.3$ seconds at $n=200,500,1000$,
scaling by ten at $R=10^{4}$), so the gap widens without bound; what the
cost law implies at $10^{4}$--$10^{6}$ alternatives is an extrapolation
and is deferred to Appendix~\ref{app:extrap}. Every timing is
regenerated by the committed \texttt{bench.py}.

Section \ref{sec:alternatives} places the method among the alternatives.
Section \ref{sec:lattice} gives the independent lattice race; Section
\ref{sec:grammar} extends it across the covariance grammar, factor to blocks
to trees, each with its recursion written out. Section \ref{sec:jac} gives
the Jacobian operator and the inversion. Section \ref{sec:general} reduces
dense covariance to the grammar family, approximately, and reports accuracy
per named random-correlation ensemble. Section \ref{sec:num} contains the
benchmarks.

\section{The alternatives}\label{sec:alternatives}

\paragraph{Deterministic quadrature.} \citet{genz1992} transformed the orthant
integral to the unit cube and integrated by adaptive and quasi-Monte Carlo rules;
the \texttt{mvtnorm} implementation is the standard in R. Accuracy is not the
issue: in the benchmark below it agrees with the lattice to the reference's own
noise. The shape is the issue. It treats one probability at a time, so the
all-$n$ vector at $n=30$ costs three hundred times the lattice, and the gap
widens with $n$.

\paragraph{Frequency simulation.} The crude frequency estimator of
\citet{lerman1981} counts argmax draws. It is unbiased for the whole vector at
once but not smooth in parameters, and its per-entry relative error explodes for
small $p_i$. Measured below: at wall clock matched to the lattice it carries
twelve to twenty times the total-variation error, records zero counts on tail
alternatives at $n=30$, and at ten times the budget remains a factor of two to
four behind.

\paragraph{GHK.} Defined above. Its dominance is documentable rather than
folklore. Stata's \texttt{cmmprobit} manual states that the likelihood
evaluator ``implements the Geweke--Hajivassiliou--Keane (GHK) algorithm to
approximate the multivariate distribution function''
\citep{statacorp2023,gates2006}, and the command's entire
\texttt{intmethod} menu (Hammersley, Halton, or pseudorandom point
sets) selects only the sequence feeding the same simulator: the
market-leading implementation offers no non-GHK method at all. R's
\texttt{mlogit} and \texttt{mvProbit} default to GHK likewise
\citep{croissant2020}, and the GHK-based \texttt{mvprobit} of
\citet{cappellari2003} carried the method to applied Stata work at large. Two properties matter
for comparison: it is smooth in parameters, which frequency counting is not,
and it prices one alternative per run. \citet{stern1992} gave a related
decomposition with an analytic leading term.

\paragraph{Bayesian data augmentation.} \citet{albert1993} and
\citet{mcculloch1994} sidestep the integral entirely by Gibbs sampling the latent
utilities; \citet{imai2005} implement the marginal data-augmentation version
in the \texttt{MNP} package, and \citet{loaizamaya2022} carry the Bayesian
route to large choice sets by variational methods. For Bayesian inference
this is a complete and often excellent answer. It is simply answering a different question than the
one benchmarked here: the cost moves into the Markov chain, the likelihood
surface is never formed, and choice probabilities appear only as posterior
functionals, so there is no pricing call to time or to check.

\paragraph{Analytic approximations.} \citet{clark1961} matched moments of
pairwise maxima; \citet{mendell1974} and \citet{solow1990} approximate by
sequential conditioning on first and second moments. These are fast, sometimes
startlingly accurate, and carry no error control, and the benchmark below
measures all three properties at once: Mendell--Elston prices the $n=10$ vector
in two milliseconds within $2\times10^{-3}$ of the reference, then drifts to
$3\times10^{-2}$ at $n=30$ with tail probabilities off by a factor of nearly
four, and nothing in the output announces the change. The Clark-type recursion
behaves alike. The family is current practice, not history:
\citet{bhat2011} builds the MACML estimator on precisely these
approximations, and the \texttt{Rprobit} package implements it
\citep{bauerRprobit}, with its asymptotics examined by
\citet{batram2019}. Expectation propagation
is the modern member of the same family. We benchmark representative
sequential moment approximations through the classical members; MACML's
composite-likelihood construction around its analytic ingredient is its own
estimator and is not itself benchmarked here.

\paragraph{Minimax tilting.} \citet{botev2017} computes single Gaussian
rectangle probabilities with bounded relative error deep into the tails, in
dimensions up to roughly a thousand. It is the accuracy reference, and we use it
as the referee for our own tail claims; it prices one probability per run.

\medskip
Every smooth, likelihood-grade method above prices one probability at a
time; frequency counting obtains the vector but not derivatives. A single
likelihood contribution is content with one probability; the tasks that
surround it are not. Share inversion, prediction across the menu, welfare,
and derivatives with respect to all $n$ utilities each want the whole
vector at every parameter value. Sharing the field across
alternatives is the algorithmic point of this paper, and it is what makes the
$n$-fold and the derivative costs disappear for those tasks.

\section{The lattice race}\label{sec:lattice}

In the min-wins convention of the introduction, let $X_i=\mu_i+\sigma_i
\epsilon_i$ with $\epsilon_i$ i.i.d.\ from a standardized base density (normal,
Gumbel, or any density supplied as code), and write $f_i$, $S_i$ for the density
and survival of $X_i$. From here on $\mu_i$ is a \emph{time}
location, the negative of the introduction's utility location: raising
it makes contestant $i$ worse and lowers $p_i$. Inversion returns time
locations, and utility coefficients are their negation.

The algorithm evaluates \eqref{eq:field-intro} on a lattice $x_1<\dots<x_L$:
\begin{enumerate}\itemsep2pt
\item compute $\log S_j(x_\ell)$ and $\log f_j(x_\ell)$ for all $j,\ell$;
\item accumulate the field $L(x_\ell)=\sum_j \log S_j(x_\ell)$;
\item for each $i$, form the integrand
      $f_i(x_\ell)\,e^{L(x_\ell)-\log S_i(x_\ell)}$, which is contestant $i$'s
      density times everyone else's survival, with $i$ divided out of the shared
      field in log space;
\item sum times the lattice spacing, and normalize the vector.
\end{enumerate}
The cost is $O(nL)$: the field is built once, not once per contestant. Log-space
arithmetic keeps hopeless contestants finite (their $e^{L-\log S_i}$ underflows
gracefully rather than dividing zero by zero), and the exponent is clipped at
$-745$.

The lattice itself is placed on the winner's bulk, not the ability span. A
hopeless contestant wins only by running a winner-class time, so only the region
where the winning time lives matters. Bisect the conservative winner c.d.f.\
envelope $G(x)=1-\prod_j S_j(x)$, built from the caller's own base survival,
to $[G^{-1}(\delta),G^{-1}(1-\delta)]$ and pad
two standard deviations, which makes the truncated tail mass negligible
against the accuracy floors reported below. Two details make that
description literally true rather than nearly true. The interval is
bracketed by geometric expansion before it is bisected: nine standard
deviations hold the $\delta$-quantiles of a Gaussian race but not of an
arbitrary polynomial tail, and bisecting an interval that does not
bracket the root converges to its endpoint while appearing to succeed.
And $\delta$ is a request rather than a guarantee, because the two error
sources trade against each other: a $t_4$ race puts its $10^{-12}$
winner quantile $1456$ standard units out, and honoring that at $257$
points would spread the lattice to a spacing of $5.7$ and lose far more
to discretization than truncation ever cost. When the requested window
will not fit the point budget at a spacing of half the tightest
performance standard deviation, $\delta$ is relaxed by factors of one
hundred until it does, and the achieved value is reported. The window is
then exactly the stated construction at the $\delta$ it names. The envelope must use the
supplied base, not a normal surrogate: a Student-$t$ base at
$\nu=2.5$ has polynomial tails that a normal envelope clips, costing
$5\times10^{-3}$ of total variation at $n=40$ against
$5\times10^{-5}$ once the true survival sets the window. Measured: $33$
points on this window reach accuracy that a $500$-point ability-span lattice
cannot, because the span window's own truncation floors near $6\times10^{-11}$;
on fields with many hopeless contestants the bulk window is one hundred times
more accurate at equal budget.

The same lattice carries the derivatives. A Jacobian evaluated on a
different grid, or one that treats the lattice total as stationary when
it is not, is the derivative of a neighbouring map rather than of the
one that was evaluated; both matter at the accuracy this paper claims.
The shipped full and row Jacobians therefore build the window through
the same routine as the forward pass, integrate the own-coordinate
derivative $\partial f_i/\partial\mu_i$ instead of imposing it from the
continuum's zero-row-sum identity, and apply the quotient rule to the
lattice normalization: with unnormalized masses $a$, total
$T=\mathbf{1}^{\top}a$ and $p=a/T$, the returned matrix is
$(A-p\,\mathbf{1}^{\top}A)/T$ with $A=\partial a/\partial\mu$. Against
central differences of the forward map this is exact to
$2\times10^{-11}$ relative on a coarse $65$-point lattice and
$1\times10^{-6}$ under a $t_4$ base, where the previous continuum
surrogate stood at $1\times10^{-8}$ and $1\times10^{-5}$; on a fine
Gaussian lattice the two agree, which is why the surrogate survived
earlier scrutiny. Row sums now vanish as a property of the quadrature
rather than by construction.

Two byproducts are cheap once the field is built. Any specified photo-finish
tie density is one more lattice sum; \S\ref{sec:jac} identifies it with the
off-diagonal Jacobian. A removal counterfactual, the probability $j$ wins
after deleting $i$, is one more division in the continuum, but the
\emph{numerical} field must be rebuilt on a window that covers the
post-removal winner: delete a dominant favorite and the surviving race
lives in the original race's second-order-statistic region, where the
winner-bulk lattice carries no mass at all. With $\mu=(-20,0,1)$ the bulk
window is about $[-29,-11]$, the post-removal truth is
$\Phi(1/\sqrt2)=0.760$, and dividing the favorite out on the original
grid leaves $3\times10^{-28}$ of unnormalized mass, whose renormalization
is numerical debris. The implementation therefore integrates removals on
their own span window and checks the unnormalized mass, raising on a
defect rather than renormalizing it away. The full $n\times n$ matrices of
such queries are $\Omega(n^2)$ outputs and are priced per query, not free.

\section{The covariance grammar}\label{sec:grammar}

Three structures make contestants conditionally independent given a few
shared variables, in increasing order of hierarchy: a global factor, factors
per cluster, and a tree of uniform shared effects. Each keeps the shared
field and the $O(nLQ)$ bill; each is exact, not fitted. The approximate
reduction of a dense covariance to this family is a separate matter and has
its own section (\S\ref{sec:general}).

\subsection{Factor}\label{sec:factor}

Let $X=\mu+Vf+\mathrm{diag}(\sqrt D)\,\epsilon$ with $f\sim N(0,I_k)$
independent of $\epsilon$, so $\Sigma=VV'+\mathrm{diag}(D)$. Conditional on $f$
the performances are independent with means $\mu+Vf$, and
\begin{equation}
p \;=\; \sum_q w_q\, p^{\mathrm{ind}}(\mu+Vf_q),
\end{equation}
a $Q$-node quadrature over runs of the independent algorithm, cost $O(nLQ)$. In
English: average the independent race over the shared luck. The node rule is adaptive rather than fixed by
rank: a tensor Gauss--Hermite grid while it stays under $10^{5}$ nodes,
so ranks three and four still use it; scrambled Sobol once the
integrand is sharp or the tensor would exceed that budget; and an
equal-weight midpoint-quantile grid at extremely sharp rank one, where
Gauss--Hermite's clustered abscissae are the wrong family for a
near-step integrand. Conditioning on the factor is
not new: \citet{butler1982} built exactly this quadrature for the one-factor
random-effects probit, and the \texttt{lpRR}/\texttt{slpRR} interfaces of
\texttt{mvtnorm} integrate over the factor dimensions of $BB'+\mathrm{diag}(D)$
by Monte Carlo. What those constructions price is one alternative per
integral. The field is the new part: one pass shares the survival product
across all $n$ alternatives, and \S\ref{sec:num} measures the difference
against the per-alternative version of this same quadrature.

One default matters. When $D_i\ll\lVert V_i\rVert^2$ the conditional race is
nearly deterministic and the integrand over $f$ is nearly a step function; a
fixed 15-node rule then silently loses up to five percent of total variation, and
the loss is invisible to lattice refinement because it is points-invariant.
The default order
now scales with the sharpness $\max_i \lVert V_i\rVert/\sqrt{D_i}$, capped by
rank.

The same pass returns the own-parameter slopes $\partial p_i/\partial\mu_i$
at no extra cost (negative in this min-wins convention: raising a time mean
lowers the win probability), which \S\ref{sec:jac} uses as the Newton
preconditioner.

\subsection{Blocks and nested}\label{sec:blocks}

Assign contestant $i$ to cluster $c(i)$ with loading $v_i\in\R^r$ and let
\begin{equation}
X_i=\mu_i+v_i'a_{c(i)}+\sqrt{D_i}\,\epsilon_i,
\qquad a_c \ \text{i.i.d.}\ N(0,I_r),
\end{equation}
a block-diagonal rank-$r$-plus-diagonal covariance: teams, conferences, sibling
products. Two independences do the work. Conditional on its own cluster's effect
a contestant is independent of its teammates; and distinct clusters are
independent outright. So the joint survival factorizes across clusters,
\begin{equation}\label{eq:blockfield}
\Pr\{X_j>x\ \text{for every } j\} \;=\; \prod_c G_c(x),
\qquad
G_c(x)\;=\;\E_a \prod_{j\in c} S_j(x\mid a)
       \;=\;\sum_q w_q\, e^{S_c(x,a_q)},
\end{equation}
with $S_c(x,a)=\sum_{j\in c}\log S_j(x\mid a)$ the cluster's conditional
log-survival. In English: $G_c(x)$ is the probability that every member of
cluster $c$ survives past $x$, with the cluster's shared luck integrated out.

Contestant $i$ then wins against a cavity field:
\begin{equation}\label{eq:cavity}
p_i \;=\; \int h_i(x)\prod_{c\ne c(i)}G_c(x)\,dx,
\qquad
h_i(x) \;=\; \sum_q w_q\, f_i(x\mid a_q)\,e^{S_{c(i)}(x,a_q)-\log S_i(x\mid a_q)}.
\end{equation}
The term $h_i$ couples $i$'s win density to its own teammates within the shared
draw of $a_{c(i)}$; the cavity product supplies the rest of the field. The cost
is $O(nLQ_a)$ whatever the number of clusters, because the per-cluster factors
are accumulated once. Rank two is special: the conditional integrand is smooth
and a tensor Gauss--Hermite rule beats scrambled Sobol at equal nodes by a wide
measured margin.

\paragraph{Nested.} Add one global factor with couplings $g_i$ and a dial
$\gamma\in[0,1]$: conditional on the global draw $f_q$ the model is a block race
at shifted abilities, so
$p=\sum_q w_q\,p_{\mathrm{block}}(\mu+\gamma g f_q)$, a finite mixture of block
races. At $\gamma=0$ the clusters are isolated, at $\gamma=1$ fully
coupled, and the covariance contribution scales as $\gamma^{2}gg'$.

\subsection{Trees}\label{sec:tree}

Let a rooted tree have the leaf clusters at its bottom and internal nodes $t$
above, and give node $t$ a scalar effect $b_t\sim N(0,1)$ applied with uniform
strength $\lambda_t$ to every leaf beneath it:
\begin{equation}
X_i=\mu_i+v_i\,a_{c(i)}+\sum_{t\,\in\,\mathrm{anc}(c(i))}\lambda_t b_t
    +\sqrt{D_i}\,\epsilon_i .
\end{equation}
The implied covariance is hierarchical: two contestants share the variance of
exactly their common ancestors. The root's effect is common to every
contestant, hence choice-irrelevant, and is fixed to zero. The uniform strength is the crucial restriction,
because a shared effect that hits every subtree member equally moves the whole
subtree's field by a translation of its argument. Messages therefore remain
functions of one variable. Per-leaf loadings on internal effects would force
two-dimensional message tables.

The algorithm is two passes over the same lattice.

\paragraph{Upward.} Each leaf cluster carries its block factor
$G_c(x)=\sum_q w_q e^{S_c(x,a_q)}$ from \eqref{eq:blockfield}. Each internal
node, visited deepest first, combines its children under its own effect:
\begin{equation}\label{eq:up}
G_t(x)\;=\;\sum_q w_q \prod_{c\,\in\,\mathrm{children}(t)} G_c\!\left(x+\lambda_t
a_q\right).
\end{equation}
In English: $G_t(x)$ is the probability that every leaf below $t$ survives past
$x$, with $t$'s effect integrated by quadrature and each child's message shifted
because the effect moves that child's whole subtree at once. Shifted evaluations
are linear interpolations on the lattice.

\paragraph{Downward.} The root's cavity is one. Each node, visited shallowest
first, receives the field of everything outside its subtree:
\begin{equation}\label{eq:down}
R_t(x)\;=\;\Bigl[\sum_q w_q\,R_{\mathrm{parent}(t)}\!\left(x-\lambda_{\mathrm{parent}(t)}
a_q\right)\Bigr]\ \prod_{s\,\in\,\mathrm{siblings}(t)} G_s(x).
\end{equation}
The sibling messages sit outside the parent-factor quadrature. This
placement is correct but not obvious: the parent effect shifts the
candidate and every sibling subtree by the same amount, so it cancels
from every candidate-versus-sibling comparison; only the cavity from
outside the parent must be shifted and averaged. Contestant $i$ then wins against its own leaf cavity,
$p_i=\int h_i(x)\,R_{c(i)}(x)\,dx$ with $h_i$ as in \eqref{eq:cavity}. The cost
is $O(nLQ)$ plus $O(L\,Q)$ per tree node.

A depth-one tree with zero strengths reproduces the block race to machine
precision, which is tested. Against a $2^{22}$-path common-random-numbers
referee the recursion sits at the simulation noise floor on every resolvable
entry at depths one through four (root-mean-square $z$ between $0.75$ and
$1.1$). Entries below the referee's resolution agree with minimax tilting
\citep{botev2017} to $0.3$ percent relative at probabilities down to
$7\times10^{-8}$, which is within the tilting estimate's own error. Given
adequate factor quadrature, tail accuracy is carried by the lattice, and the
grammar kernels price tails that simulation does not resolve.

\begin{remark}[A connection to financial practice]\label{rem:hrp}
Long-only portfolio weights could be considered as race probabilities, and
there is a tree-race covariance under which hierarchical risk parity's
bisection geometry
\citep{lopezdeprado2016} is exactly a race: give internal node $t$ strength
$\lambda_t^2=\rho_t-\rho_{\mathrm{parent}(t)}$ with $\rho_t=\max(1-2h_t^2,0)$
for merge height $h_t$, and leaf $i$ idiosyncratic variance
$1-\rho_{\mathrm{first}(i)}$. The shipped \texttt{from\_linkage} retains
the root at its cophenetic value, so its implied correlation is the
floored cophenetic matrix itself. Fixing the root effect to zero instead
gives a contrast-equivalent race, the root being choice-irrelevant: the
two agree on every pairwise difference variance, but the zero-root
model's raw correlations are $(\rho_{ij}-\rho_0)/(1-\rho_0)$ rather than
$\rho_{ij}$, which matters only to a reader reading correlations off the
model rather than choices out of it.
The floor is forced, because uniform shared effects cannot
represent negative dependence, and omitting it silently inflates other implied
correlations above one. The identity turns dendrogram-consistent concentration
limits into contest inversions, but portfolio applications are not this paper's
subject.
\end{remark}

\section{Jacobians, tie densities, inversion}\label{sec:jac}

The Jacobian is a boundary flux, and saying so once derives its formula,
its symmetry, its conservation law and its graph structure together.
Contestant $i$ wins on the polyhedral cone
$R_i=\{x: x_i\le x_k\ \forall k\}$, so $p_i(\mu)=\int_{R_i}q_\mu$ with
$q_\mu(x)=q_0(x-\mu)$; the cones meet along the photo-finish faces
$F_{ij}=\{x_i=x_j\le x_k,\ k\ne i,j\}$. Because $\mu$ enters only as a
translation, $\partial_{\mu_j}q_\mu=-\partial_{x_j}q_\mu$, so for
$j\ne i$
\[
  \frac{\partial p_i}{\partial\mu_j}
   =-\int_{R_i}\frac{\partial q_\mu}{\partial x_j}\,dx
   =-\int_{\partial R_i}q_\mu\,n_j\,dS ,
\]
the second equality being the divergence theorem, the Euclidean special
case of the generalized Stokes theorem, applied to the vector field
$q_\mu e_j$. Take $q_0\in C^1$ with $q_0$ and $\nabla q_0$ integrable and
$|x|^{n-1}q_0(x)\to0$, which the Gaussian law of the theorem satisfies
outright; integrability alone would not license either the
differentiation under the integral or the flux at infinity. The
contribution over the sphere at infinity then vanishes, and the only
face of $\partial R_i$ whose outward normal has a
nonzero $j$-component is $F_{ij}$, where that normal is
$(e_i-e_j)/\sqrt2$.

\begin{proposition}[Photo-finish derivatives are a boundary flux]
\label{prop:tie}
For $j\ne i$,
\begin{equation}\label{eq:flux}
\frac{\partial p_i}{\partial\mu_j}
 \;=\;\frac{1}{\sqrt2}\int_{F_{ij}}q_\mu\,dS\;\ge\;0 ,
\end{equation}
the probability flux across the surface on which $i$ and $j$ dead-heat
ahead of the field. Parametrizing that face by the common value
$x_i=x_j=t$ contributes a surface factor $\sqrt2$ which cancels, giving
\begin{equation}\label{eq:tiecorr}
\frac{\partial p_i}{\partial\mu_j}
 =\int \varphi_{ij}(t,t)\,
   \Pr\{X_k>t\ \forall k\ne i,j \mid X_i=X_j=t\}\,dt ,
\end{equation}
and under independence
\begin{equation}\label{eq:tie}
\frac{\partial p_i}{\partial \mu_j}
  \;=\; \int f_i(x)\,f_j(x)\prod_{k\ne i,j}S_k(x)\,dx .
\end{equation}
Under conditional independence the same factorization holds node by node
and is averaged over the shared variables. Each row of the Jacobian sums
to zero.
\end{proposition}

Everything else follows from the geometry. Symmetry is $F_{ij}=F_{ji}$:
one surface, one number, so $w_{ij}=\partial p_i/\partial\mu_j$ is
symmetric and nonnegative. The zero row sum is translation invariance, a
conservation law: what leaves one cone enters its neighbours. Writing $B$
for the oriented incidence matrix of the photo-finish graph and $W_e$ for
the diagonal of edge conductances $w_{ij}$,
\[
   J=-B^{\top}W_eB,
   \qquad
   (Jh)_i=\sum_{j\ne i}w_{ij}\,(h_j-h_i),
\]
the discrete analogue of $-\operatorname{div}(a\nabla h)$: $Bh$ takes
differences across photo-finish boundaries, $W_eBh$ is the flux across
them, $B^{\top}W_eBh$ the net flow into each winning region. The
mass-one identity is the same theorem in dimension one, since
$\sum_i f_i\prod_{j\ne i}S_j = \tfrac{d}{dt}(1-\prod_j S_j)$
integrates to one; this is why the implementation's unnormalized mass
checks are diagnostics rather than cosmetics, a material defect meaning
the lattice missed the region, not that normalization was needed.

The Jacobian is applied as an operator:
individual entries and Jacobian--vector products are byproducts of the same
field pass (per quadrature node a Gram structure over the lattice), while
materializing all $n^{2}$ entries is a choice with $\Omega(n^{2})$ cost that
nothing downstream requires. Three derivative objects should be
distinguished: the continuum identity \eqref{eq:tie}; its quadrature
evaluation; and the exact derivative of the normalized, adaptive-lattice
numerical map, which differs from the first two through the normalization
quotient and grid motion. The implementation's grid form differentiates the rectangle sum itself
and applies the normalization quotient
$(v-p\,\mathbf{1}^{\top}v)/T$, with $T$ the unnormalized total. That
quotient is not optional: without it the returned directional derivative
differs from central finite differences of the returned map by
$2\times10^{-3}$ at $L=25$. With it, agreement is $3\times10^{-11}$
from $L=101$ upward, the production range. A residual $6\times10^{-4}$
survives at $L=25$ because the window itself moves with $\mu$; that is a
grid-motion term, which no fixed-grid derivative contains and which a
finite difference recomputing its own window will always see. The
continuum form, whose symmetry serves preconditioning, is exposed
alongside. The stable pair factorization is
$[f_i e^{L-\log S_i}]\times[f_j/S_j]$: the first factor is bounded, and the
hazard $f_j/S_j$ grows linearly in the normal base and exponentially in the
Gumbel base, either of which the exponent clip absorbs. The symmetric
square-root split overflows where survivals vanish.

The block Jacobian is exact, with the within-cluster term computed under the
cavity; the nested Jacobian is the mixture of block Jacobians; the tree Jacobian
is exact within a cluster and Gram-approximate across clusters, which suffices
for Newton because residuals always use the exact forward map.
Feasibility-critical callers fall back to finite differences.

The inversion rests on a convex program.

\begin{theorem}[The inverse exists and is unique on contrasts]\label{thm:inverse}
Fix a covariance of full support (any member of the grammar with every
$D_i>0$) and let $W(\mu)=\E\min_i(\mu_i+\epsilon_i)$,
$\epsilon\sim N(0,\Sigma)$. Then $W$ is concave and differentiable with
$\nabla W(\mu)=p(\mu)$, and its Hessian is $-L(\mu)$, the negative of the
graph Laplacian whose edge weights are the tie densities \eqref{eq:tie}.
Every tie density is strictly positive, so $W$ is strictly concave on the
contrast space $\mathbf{1}^{\perp}$. For any target with $p^{\star}_i>0$
and $\sum_i p^{\star}_i=1$, the program
\[
\max_{\mu}\; W(\mu)-\langle p^{\star},\mu\rangle
\]
is constant along $\mathbf{1}$, strictly concave and coercive on
$\mathbf{1}^{\perp}$, and its unique mean-zero maximiser is the inverse:
$p(\mu^{\star})=p^{\star}$. If some $p^{\star}_i=0$ the supremum is
approached only as the contrast
$\mu_i-\min_{j:\,p^{\star}_j>0}\mu_j\to+\infty$ (the raw coordinate
statement depends on the gauge); no finite inverse exists, and a
floored target recovers a lower bound on that contrast.
\end{theorem}

\begin{proof}
A minimum of affine functions of $\mu$ is concave and expectation
preserves this. Ties have measure zero under full support, so the
minimiser is almost surely unique and Danskin's theorem gives
$\partial W/\partial\mu_i=\Pr\{i\ \text{wins}\}$.

The Hessian is Proposition~\ref{prop:tie}, which assumed no
independence: its off-diagonals are the boundary fluxes
\eqref{eq:tiecorr} under the full correlated law, and its rows sum to
zero by the same translation invariance $W(\mu+c\mathbf{1})=W(\mu)+c$
that fixes the diagonal without further computation. The Hessian is
therefore $-L(\mu)$ with $L$ the weighted graph Laplacian of the
photo-finish graph; under full support ($D_i>0$, so every pairwise difference has
positive variance) each edge weight is strictly positive, the tie graph
is complete, and $L\succ0$ on $\mathbf{1}^{\perp}$: $W$ is strictly
concave there.

Coercivity follows from the elementary bound
$W(\mu)\le\min_i\mu_i$: for each $i$,
$\min_k(\mu_k+\epsilon_k)\le\mu_i+\epsilon_i$, and $\E\epsilon_i=0$.
For a nonconstant unit contrast $d$ and interior target,
$\langle p^{\star},d\rangle>\min_i d_i$, because $p^{\star}$ puts
positive mass on a coordinate where $d$ exceeds its minimum; hence
\[
W(td)-t\langle p^{\star},d\rangle
 \;\le\; t\bigl(\min_i d_i-\langle p^{\star},d\rangle\bigr)
 \;\longrightarrow\;-\infty
\]
linearly in $t$. The decay is uniform, not merely raywise:
$d\mapsto\langle p^{\star},d\rangle-\min_i d_i$ is continuous and
strictly positive on the compact unit sphere of $\mathbf{1}^{\perp}$,
so it attains a positive minimum $\kappa$, and the displayed inequality
at $t=\lVert\mu\rVert$ gives $F(\mu)\le-\kappa\lVert\mu\rVert$ for every
$\mu\in\mathbf{1}^{\perp}$, writing $F$ for the objective --- which is
coercivity for arbitrary diverging sequences. (The bound is on $F$
itself, not on $F$ relative to $F(0)$: the objective may climb toward
its maximizer before the linear decay takes over.) Constancy along
$\mathbf{1}$ is translation invariance
again, exactly when the target sums to one. A strictly concave coercive
function on $\mathbf{1}^{\perp}$ attains a unique maximum, where the
gradient vanishes: $p(\mu^{\star})=p^{\star}$.
\end{proof}

The program is the Legendre dual of the social surplus function of
\citet{mcfadden1981}, and share inversion as a convex problem is
established territory: \citet{berry1994} inverts market shares in
logit-family demand, \citet{galichon2022} identify the general random-utility
inversion with optimal transport, and \citet{chiong2016} compute it by
linear programming in the discrete case. Closest to an algorithm is
\citet{li2018}, who for arbitrary non-atomic bases writes the inversion
as an unconstrained convex minimization whose objective gradient is the
share map, with globally convergent superlinear methods --- the
formulation and the optimizer, without a fast evaluator of the shares
and their Jacobian at large $n$. The shared field is that evaluator, so
the two are complements. What the field representation adds
is not the convex program but its derivatives: the Hessian is the
tie-density Laplacian of the same field pass, so gauge-fixed Newton runs
matrix-free at $O(nLQ)$ per iteration.

Identification lives on the same contrast space. Choice probabilities are
invariant to a common utility shock, so the model depends on $\Sigma$ only
through $P\Sigma P$ with $P=I-\mathbf{1}\mathbf{1}^{\top}/n$: a factor
proportional to $\mathbf{1}$, or a tree-root effect applied to every leaf,
is choice-irrelevant. The right fitting objective for a dense covariance
is therefore the projected residual,
\[
\min_{V,\,D\ge0}\;\bigl\lVert P\,(\Sigma-VV'-\mathrm{diag}\,D)\,P\bigr\rVert_F^2,
\]
not a fit of $VV'+\mathrm{diag}(D)$ \emph{to} $P\Sigma P$ (the projection
of a rank-plus-diagonal matrix is no longer rank-plus-diagonal, so the
second problem manufactures error the first does not have). The reference
implementation minimizes the stated objective
(\texttt{factor\_model\_projected}: alternate a top-$k$ PSD step for $V$
in the contrast basis with the exact nonnegative diagonal solve, whose
normal equations are $(P\!\circ\!P)\,d=\mathrm{diag}(P(\Sigma-VV')P)$),
and \S\ref{sec:general}'s pipeline consumes only projected residuals
downstream, so an input already of the form $VV'+\mathrm{diag}(D)$ is
returned undistorted. The introduction's identified-class sentence is to
be read in this sense: structured covariance is the estimable and
computable restriction on $\mathbf{1}^{\perp}$, not a claim that
structure exhausts identification.

The shipped iteration is stated so it can be checked against the code
(\texttt{abilities\_from\_race}).

\begin{center}
\fbox{\parbox{0.92\linewidth}{%
\textbf{Inversion} (all grammars; min-wins). Given a target $p^{\star}$
with positive entries summing to one, set $t=\log p^{\star}$ and
initialize $\mu^{0}=-(t-\bar t)/2$. Repeat until
$\max_i|r_i|<10^{-8}$:
\begin{enumerate}\setlength{\itemsep}{0pt}
\item one field pass: $\hat p=p(\mu)$, and own-slopes
$s_i=\partial p_i/\partial\mu_i$ (factor grammar: the same pass;
block/nested/tree: the slopes of the independent race at matched total
variances, one extra $O(nL)$ pass);
\item log residual $r_i=\log\hat p_i-t_i$; log-domain slope
$d_i=\min(s_i/\hat p_i,\,-10^{-6})$;
\item damped diagonal Newton step
$\mu_i\leftarrow\mu_i-\operatorname{clip}(\alpha\,r_i/d_i,\pm2)$,
then recenter $\mu\leftarrow\mu-\bar\mu$.
\end{enumerate}
Damping: $\alpha=1$ ($0.7$ at $n=2$, where the undamped update on the
$K_2$ photo-finish graph is a two-cycle); for the structure grammars
$\alpha=0.7$, halved whenever the residual fails to shrink.}}
\end{center}

This is a diagonally preconditioned Newton iteration on log
probabilities; Theorem~\ref{thm:inverse} makes the target a well-posed
problem, and the clip and recentering confine the iterate to the
contrast space. The theorem establishes well-posedness, not
convergence of this clipped iteration; the iteration is an empirically
safeguarded solver, and a convex or Newton--Krylov fallback exists in
principle should a case demand it. Round trips through the shipped front door at $n=400$
(the committed \texttt{bench.py invert}; twenty clusters, tree of depth
three): independent $0.04$\,s, tree $0.37$\,s, blocks $0.67$\,s,
factor rank-2 $1.7$\,s, nested $3.7$\,s, each to a maximum
log-probability residual below $10^{-8}$.
A zero or negative target raises rather than being silently repaired:
by Theorem~\ref{thm:inverse} it has no finite inverse. Flooring is
opt-in, and the result is then a one-sided bound on the floored
contrasts in the direction the theorem states; the returned diagnostics
name which entries were floored, whether the iteration converged and
the residual reached, and non-convergence warns rather than returning a
best effort silently. The contract is the same for every grammar: the
target is validated once, before the call splits by structure, and all
five paths return through one exit, so a block or tree inversion
reports floored entries, iteration count and convergence exactly as the
factor one does. Ignoring known structure at inversion time is
not a small error: inverting block-generated probabilities under an assumed
independent model mislocates abilities by up to three tenths of the field's
spread (the pinned configuration: $n=200$, thirty clusters, within-cluster
correlation $0.65$; maximum mislocation $2.49$ against a spread of $8.24$,
median $0.34$).

\section{Dense covariance}\label{sec:general}

The fitting procedure, stated so it can be reproduced. Write
$s_i=\sqrt{\Sigma_{ii}}$, $S=\mathrm{diag}(s)$ and
$C=S^{-1}\Sigma S^{-1}$ throughout: every fit, residual and closing
solve below is in $\Sigma$, and $C$ appears only in the clustering
distance of stage~3. The tempting alternative --- standardize to a
correlation matrix, fit there, restore scales --- fails: with
$\Sigma=SCS$, $S=\mathrm{diag}(s)$, the choice-relevant residual after
restoration is $PS(C-M)SP$, not $P(C-M)P$, and $P$ does not commute
with an unequal $S$. At $n=2$ with equicorrelated $C$ ($\rho=0.9$) and
$s=(1,2)$, the correlation-space quotient objective declares the fit
$M=(1-\rho)I$ exact while the restored model gets the only identified
quantity, the difference variance, wrong by nearly a factor of three
($0.5$ against $1.4$); the choice-irrelevant direction in correlation
coordinates is $S^{-1}\mathbf{1}$, not $\mathbf{1}$. A single global
rescaling for conditioning is harmless. The test suite pins the
two-alternative example and an exact-grammar round trip with variances
spanning two decades.

The stages, in order:
\begin{enumerate}
\item Global factors: the quotient fit of \S\ref{sec:jac}
(\texttt{factor\_model\_projected}), $k=3$ throughout, returning $(V,D_0)$
by alternation. Each half-step is optimal in its own block, so in exact
arithmetic the identified objective descends monotonically to a
coordinatewise-stationary point. Above $n=800$ the $V$-step is subspace
iteration rather than a full eigendecomposition and its subspace is only
approximate, so descent is enforced instead of assumed: the objective is
evaluated each sweep at $O(n^2)$, a sweep that fails to improve it is
discarded, and the alternation returns the best iterate. The objective is nonconvex in the pair, so no global certificate is
claimed, and the residual report exposes the value reached. Multistart
dispersion is the natural diagnostic, but it has to be measured in the
right coordinates, which is worth stating because the obvious choice is
the wrong one. Across four ensembles at $n=300$ with eight random
diagonal starts spanning two decades, every start reaches the same
objective to $4\times10^{-16}$ relative and no random start ever beats
the shipped one. On three of the four the fits agree in choices as well,
to a total variation of $10^{-14}$. On block equicorrelation they do
not: the starts agree on the objective and on $D$ to every digit, and
still price races that differ by a total variation of $0.25$. The cause
is degeneracy rather than optimization. Centering the six-block matrix
leaves a five-fold tied leading eigenvalue, so a rank-three fit selects
an arbitrary three-dimensional subspace of a five-dimensional one; every
selection is exactly optimal and they imply different races. Raising the
rank to the multiplicity makes the fit exact and collapses the
disagreement to $8\times10^{-4}$. So the diagnostic to run is dispersion
of the priced race across starts, not dispersion of the objective, which
is blind to precisely the case that matters; and when it fires, the
remedy is rank rather than a better optimizer.
\item The working residual is then
projected once, $R=P\,(\Sigma-VV'-\mathrm{diag}\,D_0)\,P$: the raw residual
still contains the choice-irrelevant common component the quotient fit
rightly ignored, and letting later stages chase it manufactures projected
error (measurably: the raw-residual variant distorts an input already of the
form $VV'+\mathrm{diag}(D)$ by $1.7\times10^{-2}$ of total variation;
the projected pipeline returns it at the $2\times10^{-4}$ node-noise
floor).
\item Blocks: average-linkage clustering on
$\sqrt{(1-C)/2}$ cut to $20$ clusters; for each cluster, the leading
eigenvector of the within-cluster block of $R$ with its diagonal zeroed,
one rank-one loading per cluster (clusters are disjoint).
\item Residual promotion: the top $m=5$ eigencomponents of what remains of $R$, diagonal
zeroed, appended as further factor columns; negative eigenvalues are not
promoted; numerically dead columns are dropped.
\item The closing diagonal
solves the identified problem's own normal equations,
$(P\!\circ\!P)\,d=\mathrm{diag}(P(\Sigma-V_{\mathrm{all}}V_{\mathrm{all}}')P)$
subject to $d\ge\ell$, with $\ell$ a small multiple of the mean
variance. The Gram is exactly
$(1-\tfrac2n)I+\tfrac1{n^2}\mathbf{1}\mathbf{1}^{\top}$, so the
constrained solve is closed form in $O(n)$: substituting
$d=\ell\mathbf{1}+x$ with $x\ge0$ leaves the same Gram with a shifted
right-hand side, and the KKT conditions make the active set a threshold
set, solved by water-filling. It is the constrained minimizer, not an
unconstrained solve with a floor applied afterwards, which the Gram's
off-diagonal coupling would leave suboptimal. At $n=2$ the Gram is rank
one, only $d_1+d_2$ is identified, and the symmetric representative is
returned.
\item A second candidate --- a pure eigenfit at the same total rank, closed by the same
diagonal solve --- is computed, and whichever candidate leaves the
smaller choice-relevant residual is kept, the pipeline winning ties:
the greedy factor-plus-blocks allocation is the wrong shape for
globally smooth covariance, and the kernel battery below measures the repair.
\end{enumerate}

The constants $k$, the cluster count and $m$ are
fixed, not tuned per ensemble; a spectral rule for $m$ was tested and
rejected (it saturates its cap for marginal gains and does not repair the
failure case below). What results is a factor-plus-blocks race priced in
one deterministic call, and the pipeline is the shipped one-call
(\texttt{fit\_covariance}; \texttt{race\_probabilities} accepts
\texttt{cov=} directly), not a research script.

The shipped call also
reports its own quality and warns on two distinct failure modes: a
large choice-relevant residual (the fit could not hold the matrix), and
a well-fitted matrix whose pricing needs a high-rank, sharp factor
integral (the quadrature, not the fit, is then the binding constraint;
both regimes appear in the kernel battery below).

At $n=2000$ the
one-call fit-and-price takes four seconds ($0.6$ of them the fit)
against thirty-six for a $10^6$-draw simulation and agrees with it at
its noise floor on the resolved bulk. For the $82$ percent of tail
alternatives with zero simulation wins, agreement in that referee is
unverifiable entry by entry, and the set is selected by the same
simulation that would have to certify it, so a fixed-set binomial bound
does not apply. Sample splitting does apply, and we run it. Fix the
zero-win set on one half of the draws, then count the wins an
independent second half places anywhere in that set: because the set is
fixed before the second half is looked at, its count is binomial in the
set's true mass and a Clopper--Pearson interval is valid. On a dense
$n=2000$ correlation with three global factors and a forty-block layer,
priced in a wide field where $63.6$ percent of alternatives take no win
in the first million paths, the model puts $1.13\times10^{-4}$ of mass
on that set and the second million paths place $111$ wins in it, a $95$
percent interval of $[9.13\times10^{-5},1.34\times10^{-4}]$. The model's
mass is inside it. Scored on the second half alone, so that no statistic
is evaluated on the draws that selected its own target, full-vector
total variation is $4.4\times10^{-3}$ against a split-half noise floor of
$5.4\times10^{-3}$: at the referee's own resolution. Tail exactness at
small $n$ still comes from the Botev referee, which prices individual
orthant probabilities rather than counting wins.

Fitting $VV'+\mathrm{diag}(D)$ \emph{to} $P\Sigma P$ is the wrong
objective (\S\ref{sec:jac}). With the identified objective and projected residuals throughout, the ensemble
study below ran a raw top-$k$ eigenfit pipeline and the identified
pipeline over twenty seeds per ensemble. On thirteen of fifteen named
ensembles the arms are statistically indistinguishable and none favors
the raw arm. On block equicorrelation the identified arm is eight
times better at the median ($2.5\times10^{-3}$ against
$2.0\times10^{-2}$), with a worst seed of $3.2\times10^{-3}$ against
the raw arm's $0.162$: the raw eigenfit spends its rank on a
near-common component, as the identification argument predicts.

Accuracy depends on the generating ensemble, so it is reported per named
ensemble of the \texttt{randomcov} library (pinned at commit
\texttt{0d27a51}), now over \emph{twenty seeds} per ensemble at $n=300$,
each seed refereed by its own $10^6$-path simulation (the committed
\texttt{run\_ensembles4.py} and \texttt{run\_kernel4.py} regenerate
everything). The TV columns are full-vector total variation against
the empirical frequency vector,
zeros included, so the referee's own counting noise and any unresolved
tail mass contribute to the reported number. That makes the statistic
upward biased in expectation, by convexity, and not a
realization-wise upper bound: on a single seed sampling error can
cancel part of the model error;
only the per-entry column is restricted, to entries with at least $25$
referee wins. TV is reported as median, ninth decile and worst over
seeds; the last column is the median per-entry absolute error on those
resolved entries.
Representative rows, identified pipeline:

\begin{center}\small
\begin{tabular}{lrrrr}
\toprule
ensemble & med TV & q90 TV & worst TV & med $|p-\hat p|$\\
\midrule
block equicorrelation & $2.5\times10^{-3}$ & $3.0\times10^{-3}$ & $3.2\times10^{-3}$ & $1.4\times10^{-5}$\\
factor $+$ sparse links & $2.9\times10^{-3}$ & $6.7\times10^{-3}$ & $1.6\times10^{-2}$ & $1.9\times10^{-5}$\\
sparse-precision graphical & $4.5\times10^{-3}$ & $5.6\times10^{-3}$ & $6.0\times10^{-3}$ & $1.5\times10^{-5}$\\
Wishart & $1.2\times10^{-2}$ & $1.4\times10^{-2}$ & $1.5\times10^{-2}$ & $2.6\times10^{-5}$\\
AR(1) & $4.6\times10^{-2}$ & $8.4\times10^{-2}$ & $1.5\times10^{-1}$ & $7.0\times10^{-5}$\\
residuals (dense-strong) & $6.8\times10^{-2}$ & $8.6\times10^{-2}$ & $9.2\times10^{-2}$ & $2.5\times10^{-4}$\\
\bottomrule
\end{tabular}
\end{center}

The medians are stable in the seed: any single seed reproduces them
to within its own noise. Eight further named
ensembles (LKJ, onion, vine, elliptope walk, animals, Marchenko--Pastur
and spiked spectra, hierarchical) fall between the third and fifth
rows; Archakov--Hansen sits with AR(1) at $4.9\times10^{-2}$.

The kernel family is the failure case, and the stratified battery
(kernel type $\times$ length scale $\times$ budget, twenty seeds per
cell, unit-square inputs) separates two distinct mechanisms that a
single ``kernel'' row had conflated:

\begin{center}\small
\begin{tabular}{llrrr}
\toprule
kernel & length scale & med TV, $m{=}5$ & $m{=}12$ & $m{=}5$, $2^{14}$ nodes\\
\midrule
RBF & $0.08$ & $3.4\times10^{-2}$ & $3.1\times10^{-2}$ & $3.3\times10^{-2}$\\
RBF & $0.2$  & $2.6\times10^{-2}$ & $2.6\times10^{-2}$ & $1.2\times10^{-2}$\\
RBF & $0.4$  & $1.6\times10^{-2}$ & $1.7\times10^{-2}$ & $8.2\times10^{-3}$\\
Mat\'ern-$\tfrac32$ & $0.08$ & $2.9\times10^{-2}$ & $2.5\times10^{-2}$ & $2.4\times10^{-2}$\\
Mat\'ern-$\tfrac32$ & $0.2$  & $2.5\times10^{-2}$ & $2.3\times10^{-2}$ & $1.6\times10^{-2}$\\
Mat\'ern-$\tfrac32$ & $0.4$  & $2.1\times10^{-2}$ & $1.7\times10^{-2}$ & $1.2\times10^{-2}$\\
\bottomrule
\end{tabular}
\end{center}

At short length scale the bottleneck is \emph{representation}:
correlation is locality, the rank needed grows with the number of
correlation lengths in the domain, and accordingly extra rank helps
($3.4\to3.1\times10^{-2}$) while extra quadrature does not. At long
length scale the fit is not the problem at all --- a rank-$27$ eigenfit
holds the RBF draw at length scale $0.4$ to a projected residual of
$10^{-3}$ --- but pricing a well-fitted smooth kernel means a
twenty-seven--dimensional \emph{sharp} factor integral (the closing
diagonal is nearly zero, so conditional races are nearly deterministic),
and there the node budget is the binding constraint: quadrupling rank
does nothing while $2^{14}$ nodes halve the error, and $2^{16}$ nodes
reach $6\times10^{-3}$ on the diagnostic draw. Stage~6 of the pipeline is what reaches this regime at all: the greedy
factor-plus-blocks allocation leaves a projected residual of $0.14$ on
the same draw that the matched-rank eigenfit holds to $10^{-3}$, a
misallocation worth up to $7\times10^{-2}$ of total variation and not a
limit of the grammar. The dispatch conclusion stands: genuinely local
covariance wants sequential conditioning with cheap sparse-precision
draws --- Mat\'ern is exactly the family with the SPDE route
\citep{lindgren2011} --- near-singular smooth covariance wants either a
larger node budget or plain simulation, and a production front door
should dispatch on the structure it detects; the shipped call's
warnings distinguish the two regimes.

\section{Benchmarks}\label{sec:num}

\paragraph{The structure-aware per-alternative comparator.} The fairest
rival is not GHK but this paper's own conditioning run the literature's
way: for each alternative separately, integrate the conditional
independent-race integrand over the same Gauss--Hermite factor nodes
(the construction of \citet{butler1982}; \texttt{mvtnorm}'s
\texttt{lpRR}/\texttt{slpRR} are its Monte Carlo form). With the same
nodes and the same lattice the two computations are algebraically
identical, and the measurement
confirms it: at $n=200$, rank~2, the per-alternative
version agrees with the shared field to total variation
$3\times10^{-15}$ and costs $601$ times as much ($22$\,s against
$36$\,ms; the committed \texttt{bench.py bm}). The comparison isolates
the paper's contribution exactly: not the conditioning, which is forty
years old, but the field shared across alternatives.


\paragraph{Against the field.} All-$n$ choice probabilities under the rank-2
factor covariance family of \S\ref{sec:factor}, against a
$2^{20}$-point QMC reference. Tail is the maximum absolute log-error over
reference probabilities in $[10^{-4},10^{-3}]$; the $n=10$ draw has none.
Frequency simulation is run at wall clock matched to the lattice and at ten
times that budget; ``zeros'' counts tail alternatives that received no draws.

\begin{center}\small
\begin{tabular}{lrrrrr}
\toprule
& \multicolumn{2}{c}{$n=10$} & \multicolumn{3}{c}{$n=30$}\\
method & time & TV & time & TV & tail\\
\midrule
lattice race & 2.5 ms & $1.8\times10^{-4}$ & 4.7 ms & $8.3\times10^{-4}$ & $0.07$\\
Genz, full vector & 46 ms & $1.8\times10^{-4}$ & 1.40 s & $8.3\times10^{-4}$ & $0.07$\\
frequency, matched time & --- & $3.8\times10^{-3}$ & --- & $1.0\times10^{-2}$ & $0.45$, 2 zeros\\
frequency, $10\times$ time & --- & $9.2\times10^{-4}$ & --- & $2.0\times10^{-3}$ & $0.21$\\
Mendell--Elston & 2.0 ms & $1.8\times10^{-3}$ & 18 ms & $3.2\times10^{-2}$ & $1.30$\\
Clark-type & 1.4 ms & $1.1\times10^{-2}$ & 17 ms & $1.6\times10^{-2}$ & $1.33$\\
\bottomrule
\end{tabular}
\end{center}

Genz and the lattice agree to the reference's noise and to each other.
Of the $300\times$ cost ratio at $n=30$, the one-at-a-time interface
supplies an unavoidable $n$-fold component; the remainder reflects the
differing integral representations, tolerances and implementations, and
we do not apportion it further. At the paper's headline scale the
$n$-fold component alone is decisive: one probability at $n=200$ costs
$0.43$ s at \texttt{maxpts} $=2.5\times10^{5}$ (agreement with the
lattice to four to five digits; six digits needs multi-second budgets
--- the committed \texttt{bench.py genz200}), so the full vector costs
$86$ s against $26$ ms for the lattice, a factor of
$3\times10^{3}$ before any accuracy matching. The sequential-conditioning family is the only entry whose error
grows silently: its tail entries are off by $e^{1.3}\approx3.7$ with nothing in
the output to say so. The race's tail figure is the reference's own noise;
lattice self-convergence places its error near $10^{-8}$.

\paragraph{Referee agreement.} A replay harness maps each benchmark
probability to its difference orthant and prices it through the two
field-standard R implementations. \texttt{mvtnorm}'s Genz--Bretz agrees with
the lattice to within its own reported error bound in every family:
$3.8\times10^{-7}$ against a bound of $8.1\times10^{-7}$ at $n=10$, and
$2.7\times10^{-8}$ against $5.8\times10^{-8}$ on a case whose smallest
probability is $3\times10^{-10}$. Botev's minimax tilting agrees to its own
Monte Carlo noise, which in relative terms is $2.2\times10^{-4}$ at
$p\approx3\times10^{-10}$: the tail claim rests on the named accuracy
reference, not only on self-convergence. On the independent family an
adaptive-quadrature referee reaches machine precision and the lattice matches
it to $1.9\times10^{-15}$. An invariance battery (two-runner closed form,
translation, permutation, symmetry) holds to $10^{-16}$ throughout, except
the two-runner comparison at $7\times10^{-9}$, which is the lattice
discretization itself.

\paragraph{Stress boundaries.} An adversarial battery pushes each failure
axis until something gives. Depth: against pure-relative adaptive
quadrature, relative error is $4\times10^{-13}$ at $p\approx10^{-30}$ and
$2\times10^{-8}$ at $p\approx10^{-50}$, with the first genuine breakdown
near $p\approx10^{-119}$, a laggard twenty-five standard deviations behind.
Inversion: the round trip $p\to\mu\to p$ holds $10^{-9}$ in log even when
the target contains a $10^{-10}$ entry or ties at the $10^{-9}$ level. The
Gumbel base reproduces softmax to $10^{-15}$ at $n=1000$; exact duplicates
split exactly, and an $\epsilon$ perturbation moves the split linearly. The
one operative boundary is factor strength. When loadings make implied
correlations exceed roughly $0.9$, the argmax integrand loses smoothness in
factor space and Gauss--Hermite converges slowly in the node count: at
loading scale $4$ the $25$-node rule still sits at $8\times10^{-3}$ total
variation. Scrambled-Sobol factor nodes restore reference-level accuracy
($4\times10^{-4}$) at the same call: heavy correlation wants
quasi-Monte-Carlo nodes, not deeper polynomial rules. The default now
escalates the family automatically past sharpness $3$, in both the Python
and R implementations (scrambled Sobol and Halton respectively), which agree
to $5\times10^{-4}$ on the worst case; the figures above are the
pre-escalation diagnosis.

\paragraph{Against GHK.} All-$n$ choice probabilities under a rank-2 factor
covariance, against a $2^{20}$-point QMC reference. TV is a bulk metric and
blind to the tails, so the same tail column as above is reported: the
maximum absolute log-error over reference probabilities in
$[10^{-4},10^{-3}]$, which for probabilities this small is the log-odds
error, and which the reference's own counting noise floors near $0.1$.

\begin{center}\small
\setlength{\tabcolsep}{4pt}
\begin{tabular}{rrrrrrrrrr}
\toprule
& \multicolumn{3}{c}{lattice race} & \multicolumn{3}{c}{GHK, $R=10^3$} & \multicolumn{3}{c}{GHK, $R=10^4$}\\
$n$ & time & TV & tail & time & TV & tail & time & TV & tail\\
\midrule
10  & 2.3 ms & $1.8\times10^{-4}$ & --- & 1.0 ms & $8.4\times10^{-3}$ & --- & 8 ms & $2.1\times10^{-3}$ & ---\\
50  & 7.0 ms & $1.7\times10^{-3}$ & $0.11$ & 21 ms & $4.2\times10^{-2}$ & $1.10$ & 214 ms & $1.2\times10^{-2}$ & $0.29$\\
200 & 25 ms  & $2.4\times10^{-3}$ & $0.15$ & 436 ms & $3.2\times10^{-2}$ & $1.50$ & 4.8 s & $1.2\times10^{-2}$ & $0.35$\\
\bottomrule
\end{tabular}
\end{center}

The $n=10$ draw has no reference probability in the tail band. The
lattice's tail figures equal the reference's own noise; GHK's are two to
ten times larger in log terms even at ten times the draws, which is the
tail blindness the TV column cannot show.

Self-convergence places the lattice error near $10^{-8}$. GHK error contracts as $O(R^{-1/2})$ and its measured complete-vector cost
grows superquadratically, with exponent $2.15$ to $2.8$ across the measured
range. The lattice is
deterministic, supplies smooth matrix-free derivatives free of simulation
noise (simulated likelihoods can be differentiated, but inherit the noise),
and prices the tails. The claim: for complete-vector pricing and inversion under the factor,
block, nested and tree covariance grammars, the shared-field
construction removes the per-alternative repetition intrinsic to
standard GHK implementations and dominates the per-alternative GHK
implementation tested here in both accuracy and cost. GHK implementations
with antithetics, quasi-random sequences or reordering heuristics would
move the constants, not the $n$-fold shape. Estimation pipelines built around GHK are a
larger object than this loop, and the estimation paragraphs
below report where simulation remains competitive at face value.
GHK also retains general rectangle probabilities and locality-structured
covariance, and \citet{botev2017} remains the reference for single
extreme tail probabilities.

\paragraph{Share inversion under sampling noise.} The first estimation
exercise is deliberately the saturated case. Simulate $T=50{,}000$ choices
from one fixed menu of $n=30$ alternatives under a known rank-2 factor
covariance and estimate the mean-zero utilities. With
the menu fixed the model is saturated over the simplex interior, so the
exact-gradient maximizer is the inversion of the observed shares and the
exercise measures how faithfully each method performs that inversion
under multinomial noise.

Each count is smoothed by one half before any arm sees it, which is the
posterior-mean share vector under the multinomial Jeffreys prior,
$\mathbb{E}[p_i\mid c]=(c_i+\tfrac12)/(T+n/2)$. (It is not the Jeffreys
MAP: the $\mathrm{Dirichlet}(\tfrac12)$ posterior has exponents
$c_i-\tfrac12$, whose mode leaves empty cells on the boundary. Read as a
penalized likelihood in simplex coordinates, add-half is the
$\mathrm{Dirichlet}(\tfrac32)$ mode.) In the saturated model the estimator
is the race inverse of that smoothed vector, and the simulated arms
maximize the same add-half-smoothed criterion, so the comparison stays
like for like. This is not cosmetic. The unpenalized saturated likelihood has \emph{no} finite
maximizer when any count is zero, which Theorem~\ref{thm:inverse}
states and this design would otherwise walk into: the smallest
alternative here has $p=2.5\times10^{-5}$, so its expected count is
$1.27$ and it is empty with probability $e^{-1.27}=0.28$. Thirteen of
the forty replications do have an empty cell, against the $11.2$ that
rate predicts, and on those an unpenalized optimizer returns a
tolerance-dependent point on a boundary sequence rather than a
maximizer. Forty replications, with the arms paired by construction
since every method sees the same draw:

\begin{center}\small
\begin{tabular}{lrrr}
\toprule
estimator & RMSE (median) & RMSE (mean $\pm$ s.e.) & median fit time \\
\midrule
exact-gradient smoothed inversion & $0.0168$ & $0.0168\pm0.0006$ & $2.4$ s \\
MSL, $R=100$ & $0.0286$ & $0.0280\pm0.0006$ & $3.5$ s \\
MSL, $R=1000$ & $0.0184$ & $0.0185\pm0.0005$ & $39.9$ s \\
\bottomrule
\end{tabular}
\end{center}

\noindent The exact path is both faster and more accurate than the
$R=100$ simulator here, and beats the $R=1000$ accuracy in a sixteenth
of its time. Because the arms share every draw, the paired differences
are the sharper statement: against $R=100$ the exact arm is better by
$0.0112\pm0.0006$, nearly nineteen standard errors, winning $39$ of $40$
replications; against $R=1000$ by $0.0017\pm0.0002$, seven standard
errors, winning $33$ of $40$. The $R=1000$ margin is small in absolute
terms, which is the honest characterization: enough simulation draws do
approach the exact answer, at forty times the cost per fit, and the
remaining gap is simulation bias rather than noise, since it survives
pairing. The timing column is from an unloaded machine at the original
eight replications; the accuracy columns are from all forty and are
insensitive to load.

\paragraph{Covariate estimation.} The overidentified case behaves
differently, and we report it as we found it. Draw observation-specific
design matrices $X_t$ ($n=10$ alternatives, $d=3$ covariates,
$T=800$ observations), utilities $\mu_t=X_t\beta$, one choice per
observation, covariance known, and estimate $\beta$ by individual-level
maximum likelihood: the exact path uses lattice probabilities with the
analytic score (one Jacobian row per observation, a single field pass);
the simulated path uses this paper's own Rust GHK with common random
numbers and finite-difference gradients. Eight replications, means with
standard errors:

\begin{center}\small
\begin{tabular}{lrr}
\toprule
estimator & RMSE$(\hat\beta)$ & median fit time \\
\midrule
exact gradient MLE & $0.0363\pm0.0066$ & $41.3$ s \\
MSL, $R=100$ & $0.0380\pm0.0073$ & $5.4$ s \\
MSL, $R=1000$ & $0.0367\pm0.0067$ & $39.6$ s \\
\bottomrule
\end{tabular}
\end{center}

\noindent All three estimators reach the statistical floor: with three
parameters and eight hundred observations the simulation error averages
out, and a well-implemented simulator at $R=100$ is the fastest of the
three at this size. The exact path's time is dominated by per-observation
Python overhead (the compiled kernel accounts for under a second of the
$41$); the comparison at $n=10$ measures plumbing, not method. What the
exact likelihood buys in estimation is not this row of the table: it is
freedom from the choice of $R$, gradients without simulation noise, the
cost law as $n$ grows, and diagnostics that simulation smooths over. The
last point is not hypothetical: on the classic Fishing data, the
unrestricted-covariance MNP likelihood evaluated exactly is
boundary-seeking (still rising at loading norms near $4000$), a ridge
that GHK's simulation noise had been regularizing by accident. The
behavior is reproducible from the committed estimator
(\texttt{winning.mnprobit}), which detects the ridge and reports it as a
\texttt{boundary\_} diagnostic rather than returning the interior
point simulation would have manufactured. For interpretability of that
norm: the parameterization is rank-2 loadings with the reference
alternative's rows fixed at zero and unit idiosyncratic variances, so
global scale is pinned by the idiosyncratic normalization and the
growing norm is a genuine choice-relevant direction, not the scale
gauge; a profile likelihood along that normalized direction is the
diagnostic a full treatment would add.

\paragraph{The Genz--Bretz score.} Our quadrature comparator used
\texttt{pmvnorm}, which exposes no score, with finite-difference gradients
($31$ vector evaluations per step at $n=30$, roughly eight times the entire
exact-gradient share-inversion fit). The \texttt{mvtnorm} manual (v1.4-2)
also documents \texttt{slpmvnorm} and, for the reduced-rank covariance
$BB'+\mathrm{diag}(D)$ of this paper's benchmarks, \texttt{lpRR} and
\texttt{slpRR}: simulated log-likelihoods and score functions that
integrate over the $K$ factor dimensions. That interface is itself the
conditioning construction of \S\ref{sec:factor} with Monte Carlo in place
of quadrature and no shared field across alternatives, which makes it the
natural structure-aware comparator for a fuller estimation study.

\paragraph{GHK protocol.} The GHK figures throughout use this paper's own
Rust implementation: per-alternative sequential conditioning on the
reference-differenced covariance with a fresh Cholesky factorization per
alternative, pseudorandom draws (xoshiro family) with a per-alternative
seed offset so that parameter evaluations share common random numbers, no
antithetics, no quasi-random sequence, no variable-reordering heuristic,
parallelized across alternatives. Cost figures are complete-vector; the
large-$n$ columns of the introduction are extrapolations of the measured
law and are displayed as such.

\paragraph{Scale.} Ten million contestants price in $245$ seconds under the
block grammar (ten thousand clusters, $257$-point lattice, the streaming
field pass; the committed \texttt{bench.py tenmillion}), and the Rust
classic-calibration kernel runs $232$ times the speed of the pure-Python
implementation. Throughput at this size is a systems statement; the accuracy
evidence lives at the scales the referees above can reach.

\paragraph{Parity.} One vector file embeds the inputs and outputs of
$22$ scenarios spanning every claim above. The four language implementations of
\S\ref{sec:avail} replay them: twenty at machine precision, the
optimizer-mediated scenarios to $10^{-6}$.


\section{Availability}\label{sec:avail}
\texttt{pip install winning} is pure Python; \texttt{winning[fast]} adds the
Rust kernels as one abi3 wheel per platform. A base-R package mirrors the full
interface, and a zero-dependency browser port drives a live demonstration at
\texttt{winning.microprediction.org}.

\begin{thebibliography}{19}\small
\bibitem[Albert and Chib(1993)]{albert1993} J.~H. Albert, S.~Chib.
Bayesian analysis of binary and polychotomous response data.
\emph{JASA} 88:669--679, 1993.
\bibitem[Berry(1994)]{berry1994} S.~Berry.
Estimating discrete-choice models of product differentiation.
\emph{RAND Journal of Economics} 25(2):242--262, 1994.
\bibitem[Butler and Moffitt(1982)]{butler1982} J.~S. Butler and R.~Moffitt.
A computationally efficient quadrature procedure for the one-factor
multinomial probit model. \emph{Econometrica} 50(3):761--764, 1982.
\bibitem[Botev(2017)]{botev2017} Z.~I. Botev.
The normal law under linear restrictions: simulation and estimation via minimax
tilting. \emph{Journal of the Royal Statistical Society B} 79(1):125--148, 2017.
\bibitem[B\"orsch-Supan and Hajivassiliou(1993)]{borschsupan1993}
A.~B\"orsch-Supan, V.~Hajivassiliou. Smooth unbiased multivariate probability
simulators for maximum likelihood estimation of limited dependent variable
models. \emph{Journal of Econometrics} 58:347--368, 1993.
\bibitem[Batram and Bauer(2019)]{batram2019} M.~Batram, D.~Bauer.
On consistency of the MACML approach to discrete choice modelling.
\emph{Journal of Choice Modelling} 30:1--16, 2019.
\bibitem[Bauer(2024)]{bauerRprobit} D.~Bauer.
\emph{Rprobit: R package for the estimation of multinomial probit models}.
\texttt{github.com/dbauer72/Rprobit}.
\bibitem[Bhat(2011)]{bhat2011} C.~R. Bhat.
The maximum approximate composite marginal likelihood (MACML) estimation of
multinomial probit-based unordered response choice models.
\emph{Transportation Research Part B} 45(7):923--939, 2011.
\bibitem[Cappellari and Jenkins(2003)]{cappellari2003} L.~Cappellari,
S.~P. Jenkins. Multivariate probit regression using simulated maximum
likelihood. \emph{Stata Journal} 3(3):278--294, 2003.
\bibitem[Chiong et al.(2016)]{chiong2016} K.~Chiong, A.~Galichon, M.~Shum.
Duality in dynamic discrete-choice models.
\emph{Quantitative Economics} 7(1):83--115, 2016.
\bibitem[Clark(1961)]{clark1961} C.~E. Clark.
The greatest of a finite set of random variables.
\emph{Operations Research} 9:145--162, 1961.
\bibitem[Croissant(2020)]{croissant2020} Y.~Croissant.
Estimation of random utility models in R: the mlogit package.
\emph{Journal of Statistical Software} 95(11):1--41, 2020.
\bibitem[Cotton(2021)]{cotton2021} P.~Cotton.
Inferring relative ability from winning probability in multi-entrant contests.
\emph{SIAM Journal on Financial Mathematics} 12(1):295--317, 2021.
\bibitem[Galichon and Salani\'e(2022)]{galichon2022} A.~Galichon,
B.~Salani\'e. Cupid's invisible hand: social surplus and identification in
matching models. \emph{Review of Economic Studies} 89(5):2600--2629, 2022.
\bibitem[Gates(2006)]{gates2006} R.~Gates.
A Mata Geweke--Hajivassiliou--Keane multivariate normal simulator.
\emph{Stata Journal} 6(2):190--213, 2006.
\bibitem[Genz(1992)]{genz1992} A.~Genz.
Numerical computation of multivariate normal probabilities.
\emph{Journal of Computational and Graphical Statistics} 1:141--149, 1992.
\bibitem[Geweke(1989)]{geweke1989} J.~Geweke.
Bayesian inference in econometric models using Monte Carlo integration.
\emph{Econometrica} 57:1317--1339, 1989.
\bibitem[Hajivassiliou et al.(1996)]{hajivassiliou1996} V.~Hajivassiliou,
D.~McFadden, P.~Ruud. Simulation of multivariate normal rectangle probabilities
and their derivatives. \emph{Journal of Econometrics} 72:85--134, 1996.
\bibitem[Imai and van Dyk(2005)]{imai2005} K.~Imai, D.~A. van Dyk.
MNP: R package for fitting the multinomial probit model.
\emph{Journal of Statistical Software} 14(3):1--32, 2005.
\bibitem[Keane(1994)]{keane1994} M.~P. Keane.
A computationally practical simulation estimator for panel data.
\emph{Econometrica} 62:95--116, 1994.
\bibitem[Lerman and Manski(1981)]{lerman1981} S.~Lerman, C.~Manski.
On the use of simulated frequencies to approximate choice probabilities.
In \emph{Structural Analysis of Discrete Data}, MIT Press, 1981.
\bibitem[Li(2018)]{li2018} L.~Li. A general method for demand
inversion. arXiv:1802.04444 (2018).
\bibitem[Lindgren et al.(2011)]{lindgren2011} F.~Lindgren, H.~Rue,
J.~Lindstr\"om. An explicit link between Gaussian fields and Gaussian Markov
random fields. \emph{Journal of the Royal Statistical Society B}
73(4):423--498, 2011.
\bibitem[Loaiza-Maya and Nibbering(2022)]{loaizamaya2022} R.~Loaiza-Maya,
D.~Nibbering. Scalable Bayesian estimation in the multinomial probit
model. \emph{Journal of Business \& Economic Statistics}
40(4):1678--1690, 2022.
\bibitem[L\'opez de Prado(2016)]{lopezdeprado2016} M.~L\'opez de Prado.
Building diversified portfolios that outperform out of sample.
\emph{Journal of Portfolio Management} 42(4):59--69, 2016.
\bibitem[McFadden(1981)]{mcfadden1981} D.~McFadden.
Econometric models of probabilistic choice. In C.~Manski, D.~McFadden,
eds., \emph{Structural Analysis of Discrete Data}, MIT Press, 1981.
\bibitem[McCulloch and Rossi(1994)]{mcculloch1994} R.~McCulloch, P.~Rossi.
An exact likelihood analysis of the multinomial probit model.
\emph{Journal of Econometrics} 64:207--240, 1994.
\bibitem[Mendell and Elston(1974)]{mendell1974} N.~Mendell, R.~Elston.
Multifactorial qualitative traits: genetic analysis and prediction of
recurrence risks. \emph{Biometrics} 30:41--57, 1974.
\bibitem[Solow(1990)]{solow1990} A.~R. Solow.
A method for approximating multivariate normal orthant probabilities.
\emph{Journal of Statistical Computation and Simulation} 37:225--229, 1990.
\bibitem[Stern(1992)]{stern1992} S.~Stern.
A method for smoothing simulated moments of discrete probabilities in
multinomial probit models. \emph{Econometrica} 60:943--952, 1992.
\bibitem[StataCorp(2023)]{statacorp2023} StataCorp.
\emph{Stata Choice Models Reference Manual: cmmprobit, asmprobit}.
Stata Press, College Station, 2023.
\bibitem[Thurstone(1927)]{thurstone1927} L.~L. Thurstone.
A law of comparative judgment. \emph{Psychological Review} 34:273--286, 1927.
\bibitem[Train(2009)]{train2009} K.~Train.
\emph{Discrete Choice Methods with Simulation}. 2nd ed., Cambridge, 2009.
\end{thebibliography}
\appendix
\section{Extrapolated GHK cost brackets}\label{app:extrap}

The lattice column below is measured; the GHK column is an
extrapolation of the two measured exponents ($2.15$ and $2.8$) far
beyond any regime in which GHK has been run, and cache effects,
parallelism, Cholesky reuse or a change of regime would move it.

\begin{center}\small
\begin{tabular}{rrrr}
\toprule
$n$ & lattice, measured & GHK at $10^4$ draws, cost-law bracket & multiplier \\
\midrule
$10^4$ & $0.18$ s & $20$ hours -- $3.7$ days & $4\times10^{5}$ -- $2\times10^{6}$ \\
$10^5$ & $2.7$ s & $0.3$ -- $6.5$ years & $4\times10^{6}$ -- $8\times10^{7}$ \\
$10^6$ & $29$ s & $46$ -- $4{,}100$ years & $5\times10^{7}$ -- $4\times10^{9}$ \\
\bottomrule
\end{tabular}
\end{center}

Nobody prices a million-alternative field with GHK; the brackets say
what the measured law implies about why. The ten-million forward-field
benchmark of \S\ref{sec:num}, which is actually run, is the meaningful
systems demonstration.

\end{document}
